\chapter{Lifting Operations and Lifting Equipment Risk Assessment}
\label{chap:Lifting Operations and Lifting Equipment Risk Assessment}

%---------------------------- SECTION ----------------------------
\section{Why this assessment matters in construction}

\paragraph{Lifting operations represent one of the highest-energy and most unforgiving hazard categories in construction. When a crane drops its load, when a MEWP overturns, when a chain block fails under dynamic loading, or when a sling parts at the point of lift, the consequences are instantaneous and typically severe. The Lifting Operations and Lifting Equipment Regulations 1998 (LOLER) were enacted precisely because the free-market tolerance of repeated lifting fatalities --- crane collapses, falling loads striking workers and members of the public, MEWP overturns --- had demonstrated that voluntary controls were insufficient to manage the inherent physics of suspended loads. LOLER, in combination with PUWER 1998 and the Construction (Design and Management) Regulations 2015, provides a comprehensive regulatory framework that imposes specific duties to plan, supervise, and carry out lifting operations safely. Despite this framework, lifting-related fatalities and major injuries continue to occur on UK construction sites with a regularity that demands the lifting risk assessment be treated as a technical document of the highest order.}

\paragraph{The complexity of lifting operations in construction is frequently underestimated at the planning stage. A tower crane erection on a city-centre site involves a machine that can weigh in excess of 200 tonnes when assembled, that must be erected in a confined urban environment, that will operate over public footpaths and occupied buildings, and that will remain on site for months or years, exposed to all weather conditions. A MEWP deployed on a ground-bearing floor slab must have its ground-bearing pressure evaluated against the slab design capacity before every deployment; the failure of a concrete slab under MEWP outrigger loading has caused multiple fatalities and serious injuries. A simple chain block used to install structural steelwork operates in an environment of dynamic shock loading, potential mis-rigging, and asymmetric load distribution that can exceed its rated capacity without any visible sign of overloading. Each of these scenarios carries fatal potential and each requires a specific, competently produced lifting plan.}

\paragraph{The human factors dimension of lifting safety is equally significant. Crane operators, riggers, slinging operatives, and appointed persons conducting lifting operations work under conditions of pressure, time constraint, and physical distance from the load that create the conditions for systematic error. A tower crane operator at 50 metres cannot see the ground-level rigging of the load; a mobile crane operator erecting a precast facade panel cannot see the connection team on the third-floor slab; a MEWP operator cannot assess the ground conditions beneath the outrigger pads from the basket. The requirement of LOLER for all lifting operations to be planned by a competent person exists to ensure that the planning process --- which can be conducted at desk level, with time, with full load data, and with engineering calculations --- compensates for the inherent limitations of the real-time operational environment.}

\paragraph{The public domain presents particular lifting challenges. Tower cranes on city-centre sites operate over public footpaths, parked vehicles, and occupied properties as a matter of course. The slew arc of a mobile crane on a street works project may pass over a public bus stop or a school entrance. The lifting plan for any operation whose load path passes over or adjacent to a public area must include specific measures to exclude the public from the hazard zone, to notify the relevant highway and planning authorities, and to obtain any necessary licences or permits. The failure to close a footpath or erect adequate exclusion barriers before a crane lift over a public area is not merely a safety failure --- it is a criminal offence under the Health and Safety at Work etc.\ Act 1974 and may engage the Corporate Manslaughter and Corporate Homicide Act 2007.}

\barquote{``No lifting operation should ever begin without a competent lifting plan that addresses the load weight, the centre of gravity, the ground conditions, the weather envelope, the load path, and the emergency procedure. A crane does not forgive improvisation.''}

\example{Tower Crane Over Public Area}{A principal contractor is constructing a twelve-storey residential tower on a brownfield site adjacent to a public footpath that serves as the primary pedestrian route between a railway station and a residential estate. The tower crane is positioned to lift precast concrete stair flights of approximately 3.5 tonnes each from an articulated delivery vehicle on the public highway to a temporary storage area on the podium slab, and then from storage to final position in the stairwell. The appointed person produces a lifting plan that identifies: the maximum rated capacity of the crane at the required radius (confirmed to be 4.2 tonnes, providing a safety margin of 700\,kg above the heaviest stair flight); the ground conditions beneath the crane base (confirmed by structural engineer to be adequate for crane base loading); the public footpath closure requirement (a temporary traffic management scheme and footpath diversion signed off by the highway authority); the exclusion zone beneath the load path (established with proprietary barrier systems and maintained by a banksman at each end of the exclusion zone); and the weather envelope (lift to cease if sustained winds exceed 20\,m/s at crane top-level anemometer). The lift proceeds without incident. A subsequent audit reveals that three similar sites operated by the same contractor have tower cranes over public areas without highway authority notification, without load path exclusion zones, and without weather monitoring --- the lifted precast units at those sites weigh within 5\% of the crane's rated capacity at the operating radius.}

%---------------------------- SECTION ----------------------------
\section{Legal and professional framework}

\paragraph{The regulatory framework for lifting operations is among the most detailed and technically specific in UK health and safety law. LOLER 1998 imposes duties that go beyond risk assessment to require specific competencies, examination regimes, and planning standards. The interaction between LOLER, PUWER 1998, CDM 2015, BS 7121, and the Supply of Machinery (Safety) Regulations 2008 creates a multi-layered compliance framework that demands genuine technical competence from those managing lifting operations on construction sites.}

\paragraph{Key frameworks relevant to this chapter include: Lifting Operations and Lifting Equipment Regulations 1998 (LOLER); Provision and Use of Work Equipment Regulations 1998 (PUWER); Health and Safety at Work etc.\ Act 1974; Construction (Design and Management) Regulations 2015; BS 7121 (Code of Practice for the Safe Use of Cranes); Management of Health and Safety at Work Regulations 1999; Supply of Machinery (Safety) Regulations 2008.}

\begin{itemize}
  \bitem{Lifting Operations and Lifting Equipment Regulations 1998 (LOLER)}{The primary instrument governing lifting operations and lifting equipment. Regulation 8 requires that every lifting operation is properly planned by a competent person, appropriately supervised, and carried out in a safe manner. Regulation 9 requires that lifting equipment used for lifting persons is examined by a competent person every six months, and that all other lifting equipment is examined at least every twelve months or at intervals specified by the competent person. Regulation 10 requires that all lifting equipment is accompanied by a report of thorough examination when first put into service and after exceptional circumstances affecting its safety. LOLER defines lifting equipment broadly to include cranes, hoists, MEWPs, gin wheels, chain blocks, shackles, slings, and any other equipment used in lifting operations.}
  \bitem{Provision and Use of Work Equipment Regulations 1998 (PUWER)}{Applies to all lifting equipment as work equipment and requires that it is suitable for its intended purpose, maintained in an efficient state, and used only by trained and competent persons. PUWER Regulation 22 specifically addresses power presses and certain lifting equipment in terms of inspection requirements. PUWER supplements LOLER by requiring the broader suitability assessment of equipment before selection and use.}
  \bitem{Construction (Design and Management) Regulations 2015}{Requires that the Construction Phase Plan addresses the arrangements for managing lifting operations on the project, including the identification of areas where crane operations will take place, the management of crane erection and dismantling, and the competency requirements for appointed persons and crane operators. The pre-construction information must include any constraints on crane positioning arising from ground conditions, buried services, adjacent structures, or overhead lines.}
  \bitem{BS 7121 --- Code of Practice for the Safe Use of Cranes}{The multi-part British Standard that provides detailed technical guidance on the planning, operation, and maintenance of cranes. BS 7121 Part 1 covers general aspects; Part 2 covers inspection, testing, and examination; Part 5 covers tower cranes; Part 6 covers MEWPs. The appointed person, a specific role defined by BS 7121, is the competent person responsible for planning each lifting operation and has authority over all crane-related activities.}
  \bitem{Management of Health and Safety at Work Regulations 1999}{Requires a suitable and sufficient risk assessment of all lifting operations, the appointment of competent persons to plan and supervise those operations, and the implementation of a management system that ensures the risk controls identified in the assessment are implemented and verified.}
  \bitem{Supply of Machinery (Safety) Regulations 2008}{Requires that lifting equipment placed on the EU/UK market carries a CE/UKCA mark indicating conformity with the essential health and safety requirements of the Machinery Directive/Machinery Regulation, and is accompanied by a Declaration of Conformity. The absence of a CE/UKCA mark on lifting equipment is a primary indicator that the equipment has not been assessed for the basic design-level safety requirements.}
\end{itemize}

%---------------------------- SECTION ----------------------------
\section{Conducting the assessment using the five-step method}

\subsection{Step 1 --- Identify Hazards}
\paragraph{Lifting hazards in construction fall into several distinct categories, each requiring specific assessment. Load failure and drop hazards include: rigging failure due to damaged, incorrect-grade, or inadequately rated slings, shackles, and lifting points; overloading of the crane or hoist due to inaccurate load weight estimates or failure to account for dynamic loading factors; load swing due to off-vertical lifts, wind loading, or snatch lifts; and load instability due to poor centre-of-gravity identification, inadequate load securing, or asymmetric rigging. Equipment failure hazards include: crane structural failure due to overloading, fatigue damage, inadequate maintenance, or inadequate ground support; MEWP tipping due to ground instability, outrigger pad failure, or overloading of the basket; chain block and lever hoist failure due to overloading or misuse; and wire rope failure due to kinking, corrosion, or wear. Operational hazards include: persons struck by the suspended load or crane jib during a lift; persons struck by the crane counterweight or slewing body during a slewing operation; MEWP operator falls from the basket; persons trapped between the load and a fixed structure during an in-position lowering; and overhead electrical contact by the crane jib, boom, or lifting line.}

\subsection{Step 2 --- Decide Who or What Is Affected}
\paragraph{Persons at risk from lifting operations include: the slinging operative responsible for attaching the load; the crane operator; banksmen guiding the load into its landing position; other tradespeople working in the area beneath or adjacent to the load path; members of the public beneath an external lift over a public area; vehicle drivers delivering materials to the crane area; MEWP operators and passengers; and operatives making structural connections to lifted steelwork or precast units. Property at risk includes: existing structures adjacent to the load path that may be struck by a swinging load; overhead electrical lines that may be contacted by the crane jib; buried structures and services that may be damaged by crane outrigger or tower crane base loading; and the lifted load itself, which may be a high-value structural component or a precision mechanical item.}

\subsection{Step 3 --- Evaluate Likelihood and Consequence}
\paragraph{The consequence of a dropped load, crane collapse, or MEWP overturn is potentially catastrophic. A three-tonne precast panel falling from a height of 15 metres has a kinetic energy at impact equivalent to a small explosive charge; anyone in the load path will suffer fatal or life-changing injury. The inherent risk of a lifting operation without a competent lifting plan, without thorough examination of the equipment, and without an established exclusion zone is Very High. The effectiveness of LOLER controls in reducing inherent risk to an acceptable residual level is high when all controls are correctly implemented: a correctly planned, correctly supervised, and correctly executed lift with thoroughly examined equipment and a competent appointed person can achieve a Low residual risk rating. The assessment must evaluate not only the probability and consequence of a load drop but also the consequence of a ground-bearing failure beneath a MEWP or crane outrigger --- a consequence that may occur without any operational error by the crane team if the ground assessment has been inadequate.}

\subsection{Step 4 --- Select and Implement Controls}
\paragraph{Every lifting operation on a construction site must have a documented lifting plan produced by a competent appointed person before the lift commences. The lifting plan must include: the precise weight and centre of gravity of the load; the crane model, configuration, and rated capacity at the required radius; the ground-bearing pressure under crane outriggers or tower crane base and the structural engineer's confirmation that the bearing surface is adequate; the load path from pick point to landing point, including identification of any overhead obstruction, adjacent structure, or public area; the rigging scheme (sling type, grade, angle, and length; shackle type and rated capacity; lifting points confirmed on the structural drawings or weighed and measured on site); the exclusion zone beneath and around the load path; the weather envelope within which the lift is approved; the communication system between crane operator and slinging operative; and the emergency procedure for a load hold if the lift must be suspended mid-operation. No lifting operation should be improvised; every change to the planned lift parameters (weight, radius, ground position) requires a formal reassessment by the appointed person.}

\subsection{Step 5 --- Record and Review}
\paragraph{The lifting plan must be retained on site and presented to the crane operator, slinging operatives, banksmen, and site manager before every lift. A pre-use inspection of the crane and rigging equipment must be conducted each day before lifting commences, with any defects recorded and equipment taken out of service until repaired. The thorough examination certificate for the crane (six-monthly for personnel lifts, twelve-monthly for goods lifts) must be current and available on site. Any incident involving a load drop, rigging failure, near-miss during a lift, or adverse weather event that required the lift to be suspended must trigger an immediate review of the lifting plan and the LOLER thorough examination status of all equipment involved. RIDDOR notification must be considered immediately following any incident involving lifting equipment.}

\example{Five-Step Application}{A structural steelwork contractor is installing a 2.8-tonne roof beam at a height of 9 metres using a 35-tonne mobile crane on an occupied commercial site. Step 1 identifies: load drop from rigging failure; ground bearing failure under outrigger pads on a former basement slab; load swing due to working in the lee of an adjacent building at wind speeds variable up to 15\,m/s; persons in the load path (site office workers using an adjacent temporary building). Step 2 identifies: slinging operative; crane operator; connection team on elevated steelwork; site office occupants. Step 3 rates inherent risk as Very High. Step 4 implements: appointed person produces lifting plan confirming 2.8\,tonnes at 8-metre radius is within 78\% of rated capacity; structural engineer confirms former basement slab capacity is adequate for 12-tonne outrigger pad loading with 100mm proprietary pads; two-sling bridle rig with certified Grade 8 chain slings rated at 4.2 tonnes each at 60-degree angle; banksman-controlled exclusion zone; site office evacuated and vacated for the duration of the lift; wind speed monitor confirms 12\,m/s maximum during lift window; lift cancelled and rescheduled when 18\,m/s gusts are forecast. Residual risk assessed as Low. Step 5: pre-lift inspection of slings, shackles, and crane checks conducted and recorded; lifting plan signed by appointed person, crane operator, and site manager.}

%---------------------------- SECTION ----------------------------
\section{Applying the hierarchy of controls}

\paragraph{The hierarchy of controls for lifting operations is shaped by the physics of the hazard: the energy of a suspended load cannot be contained by PPE or administrative instruction alone. Engineering controls --- competent mechanical design of the lift, rated and examined equipment, adequate ground support --- are the primary protective layer, supported by administrative controls that ensure the engineering systems are correctly operated and maintained.}

\begin{itemize}
  \bitem{Elimination}{Eliminate the lifting operation entirely where possible: design structures and services to allow ground-level assembly and connection before installation by alternative means; specify modular construction approaches that remove the need for multiple high-risk individual lifts; use pre-fabricated mechanical and electrical risers assembled at ground level and installed as complete modules; design building envelopes that allow facade installation from internal platforms rather than external lifting.}
  \bitem{Substitution}{Substitute higher-risk lifting methods with lower-risk alternatives: use a tower crane with an established, regularly inspected base rather than a mobile crane repositioned on variable ground for each lift; use a MEWP with a certified ground-bearing capacity that can be confirmed against slab design drawings rather than relying on estimated ground strength; use an electric chain hoist on a certified monorail system in preference to a manually operated lever hoist in a confined structural bay.}
  \bitem{Engineering controls}{Competently produced and documented lifting plans for every lift; thorough examination of all lifting equipment at LOLER-prescribed intervals; rated and marked slings, shackles, and lifting accessories with valid examination records; ground-bearing capacity assessment by a structural engineer for every mobile crane or MEWP deployment; weather monitoring with specified wind speed limits for each lifting operation; load cells on the crane hoist line for operations near the crane's rated capacity.}
  \bitem{Administrative controls}{Appointed person system with a named, competent individual responsible for planning every lifting operation; pre-lift briefing for all persons involved in the lift; exclusion zone management by a dedicated banksman; no lift to commence without appointed person sign-off on the lifting plan; LOLER examination certificates to be on site and in date before lifting operations commence; crane operators and slinging operatives to hold current CPCS cards.}
  \bitem{PPE}{Hard hat for all persons in the crane and lifting area; high-visibility vest for banksmen and slinging operatives; safety boots with steel midsoles; safety harness and inertia reel for connection operatives working at height during steelwork erection; face shield and leather gauntlets for chain sling operations where pinching and abrasion hazards are present at the lift point.}
\end{itemize}

%---------------------------- SECTION ----------------------------
\section{Cadence of assessment and review triggers}

\paragraph{The cadence of lifting operation assessment and review must match the dynamic nature of the construction site environment. Crane positions change; ground conditions change; the weights and configurations of loads change as the project progresses. The lifting plan produced for the first structural steel delivery will not be valid for the final facade element lift three months later. The appointed person must be actively engaged throughout the project's lifting operations programme, not merely at the initial planning stage.}

\begin{itemize}
  \bitem{Before every lifting operation}{Each distinct lifting operation --- each separate pick-and-place cycle involving a different load, a different crane radius, or a different ground position --- requires a lifting plan or a confirmed assessment that the existing lifting plan remains valid for the new parameters. The appointed person must confirm plan validity before every shift of lifting operations commences.}
  \bitem{After any change to crane configuration or ground position}{Any change to the crane's base position, jib configuration, counterweight, or outrigger spread must trigger a re-assessment of the rated capacity at the new radius and a re-assessment of the ground-bearing capacity at the new position. Mobile cranes repositioned between lifts represent one of the highest-risk moments in lifting operations.}
  \bitem{When weather conditions approach the specified limit}{When wind speeds approach the maximum stated in the lifting plan, the appointed person must monitor conditions and make the decision to suspend lifting operations before the limit is exceeded, not after. Lifting in borderline weather conditions based on the judgement of the crane operator in the cab is a management failure.}
  \bitem{After any rigging failure, near-miss, or lifting incident}{Any failure of a sling, shackle, or lifting point during a lift --- even one that does not result in a dropped load --- must immediately suspend all lifting operations, trigger a LOLER examination of all rigging equipment in use, and require a revised lifting plan before operations resume. Near-misses in lifting operations are indicators of systemic planning or equipment failures.}
  \bitem{At each change of lifting phase}{When the nature of the lifting programme changes --- from structural steelwork to precast concrete, from precast to M\&E plant, from M\&E to facade --- the lifting plan must be reviewed and updated to reflect the different load characteristics, rigging requirements, and ground conditions of the new phase.}
\end{itemize}

%---------------------------- SECTION ----------------------------
\section{Common failure modes}

\begin{itemize}
  \bitem{Load weights estimated rather than confirmed}{The rated capacity of a crane is expressed in tonnes; the safety of a lift depends on the load weight being known with reasonable precision before the lift commences. Estimating the weight of a precast concrete unit by reference to its approximate dimensions and assumed density, without referring to the structural drawings or the manufacturer's data sheet, introduces an error that may place the crane near or over its rated capacity at the operating radius. All loads must have a confirmed weight before the lifting plan is finalised.}
  \bitem{Ground conditions not assessed for mobile crane deployment}{The deployment of a mobile crane on ground whose bearing capacity has not been assessed by a structural engineer or geotechnical specialist is one of the most common causes of crane overturn incidents. Made ground, backfilled trenches, former basements, soft alluvial soils, and waterlogged ground have all been the ground conditions in which cranes have overturned following outrigger pad failure. The lifting plan must always include a ground-bearing assessment.}
  \bitem{Lifting accessories without valid thorough examination records}{LOLER requires that lifting accessories (slings, shackles, hooks, spreader beams, and lifting frames) are thoroughly examined by a competent person at six-monthly intervals. Construction sites frequently accumulate lifting accessories of unknown provenance, with examination records that cannot be located or that have expired. Any lifting accessory without a current, traceable examination record must be taken out of service.}
  \bitem{MEWP operated on ground not assessed for outrigger loading}{MEWPs deployed on concrete slabs, temporary road surfaces, or ground that has not been assessed for the specific outrigger pad loading of the machine are at risk of outrigger pad punch-through or settlement during the lift, causing the machine to tip. The MEWP manufacturer's load chart must be cross-referenced against the structural engineer's bearing capacity assessment for every MEWP deployment.}
  \bitem{No exclusion zone established for public or adjacent workers}{Lifting operations over or adjacent to public areas, site offices, or areas where other trades are working without an established and enforced exclusion zone represent a direct breach of LOLER and the Health and Safety at Work etc.\ Act 1974. The exclusion zone must be established before the crane is configured for the lift, not as an afterthought when the load is already in the air.}
  \bitem{Appointed person absent or role poorly defined}{The appointed person role --- the competent individual responsible for planning and supervising the lifting operation --- is frequently either absent, poorly defined, or assigned to someone who does not hold the required competency. A site manager who signs the lifting plan without the technical competence to assess its adequacy is not fulfilling the appointed person function; a crane operator who accepts verbal instructions about load weights and ground conditions without a formal lifting plan is operating without the protection LOLER requires.}
\end{itemize}

\example{Failure Pattern}{A mechanical and electrical contractor installs a 1,800\,kg rooftop air handling unit using a 12-tonne pick-and-carry rough terrain forklift. The forklift is rated at 3.2 tonnes at 500\,mm load centre; the air handling unit is 1.4 metres deep, placing its centre of gravity at approximately 700\,mm from the forklift's face. At the 700\,mm load centre, the forklift's rated capacity is approximately 2.1 tonnes --- adequate for the 1.8-tonne load, but only with a margin of 300\,kg. No lifting plan has been produced. No load centre calculation has been performed. The forklift operator has not seen the manufacturer's load chart. While manoeuvring the unit across the roof slab towards the plant room entrance, the front left wheel encounters a drainage gully cover that deflects under the combined load; the forklift tilts forward and to the left, depositing the air handling unit through the roof membrane into the building interior below. No persons are in the building beneath the penetration at the time of the incident. The investigation finds that: no lifting plan was produced; the forklift load chart was not consulted; and the roof slab design had not been checked for the dynamic loading of the loaded forklift traversing the slab. Three separate planning failures converge in a single preventable incident.}

%---------------------------- CHAPTER SUMMARY ----------------------------
\newpage
\section{Chapter Summary}

\begin{table}[H]
\centering
\begin{tabularx}{\textwidth}{>{\raggedright\arraybackslash}p{3.2cm}X}
\toprule
\textbf{Assessment Element} & \textbf{Operational Requirement} \\
\midrule
Legal/Professional Basis & LOLER 1998; PUWER 1998; CDM 2015; BS 7121; Health and Safety at Work etc.\ Act 1974; Supply of Machinery (Safety) Regulations 2008. \\
Method & Apply the full five-step process and document inherent/residual risk. \\
Controls & Competent appointed person; documented lifting plan for every lift; LOLER thorough examination of all equipment in date; ground-bearing assessment for every mobile crane and MEWP deployment; established and enforced load path exclusion zones. \\
Cadence & Before every lifting operation; after any change in crane position or load parameters; when weather conditions approach the specified limit; after any incident or near-miss. \\
Assurance & Use audit, incident learning, and governance reporting to verify effectiveness. \\
\bottomrule
\end{tabularx}
\end{table}

\begin{itemize}
  \bitem{Key takeaway 1}{Every lifting operation must be planned by a competent appointed person before the lift commences; the lifting plan must address load weight, ground conditions, equipment capacity at radius, rigging scheme, load path, exclusion zones, weather envelope, and emergency procedure.}
  \bitem{Key takeaway 2}{All lifting equipment and accessories must have a current LOLER thorough examination report; equipment without a traceable, in-date examination must be taken out of service immediately, regardless of its apparent condition.}
  \bitem{Key takeaway 3}{Ground-bearing capacity beneath mobile crane outriggers and MEWP pads must be assessed by a competent structural or geotechnical engineer before deployment; estimated or assumed ground capacity is not an acceptable basis for lifting operations.}
  \bitem{Key takeaway 4}{Exclusion zones beneath and around the load path must be established before the crane is configured for the lift, and must be maintained by a dedicated banksman for the duration of every lifting operation.}
\end{itemize}

\barquote{In Chapter~\ref{chap:Plant, Vehicles and Site Transport Risk Assessment}, we turn to the management of mobile plant, site vehicles, and pedestrian-plant segregation.}
